\section{Statistical foundations and related work}
\label{sec:related}

\subsection{Exponential families and conjugate segment models}
The exponential-family representation separates a finite-dimensional sufficient statistic from the log normalizer of the sampling distribution. Diaconis and Ylvisaker characterize conjugate priors for natural exponential families and the hyperparameter conditions under which those priors are proper \citep{diaconis1979conjugate}. In a changepoint model, the same calculation is applied separately to every candidate segment. BayesBreak uses this standard conjugate theory; its contribution is to formulate the resulting segment marginal likelihoods and posterior moment numerators so that they can be reused by the same partition recursions and by the hierarchical extensions developed below. Exactness requires the stated regularity, propriety, and finiteness conditions. The method does not assert closed-form integration for arbitrary nonconjugate parameterizations.

\subsection{Product partition models and exact Bayesian recursions}
Product partition models assign a partition probability proportional to a product of segment cohesions and combine this prior with segmentwise integrated likelihoods \citep{barry1992ppm,barry1993bayesCP}. This is the principal probabilistic lineage of the present work. Fearnhead developed exact recursions and posterior simulation for a broad class of Bayesian multiple-changepoint models with conditionally independent segments \citep{fearnhead2006exact}; Hutter derived exact Bayesian regression for piecewise-constant functions, including posterior break probabilities and model-averaged regression curves \citep{hutter2006bpcr}. Dependence between adjacent segments can sometimes be accommodated by enlarging the recursion state \citep{fearnhead2011dependence}. Bayesian Blocks provides a related local-fitness and dynamic-programming construction for astronomical time series \citep{scargle2013bayesianblocks}.

BayesBreak does not claim priority for product partition models, conjugate segment integration, exact changepoint recursions, posterior changepoint probabilities, or Bayesian piecewise-constant regression. It develops a generalized hierarchical model in which these established calculations are stated with common notation across exponential-family segment models, irregular-design priors, multiple sequences, observed groups, and a finite latent-group extension. It also keeps posterior marginalization and joint MAP optimization as separate recursions rather than treating one as an approximation to the other.

\subsection{Frequentist segmentation and optimization}
Dynamic programming also underlies frequentist segmentation with additive segment costs \citep{auger1989segment,jackson2005optpart}. Penalized formulations and pruning lead to PELT and related pruned exact algorithms \citep{killick2012pelt,rigaill2010pruned}. Wild binary segmentation and SMUCE pursue different detection and uncertainty goals \citep{fryzlewicz2014wbs,frick2014smuce}. Total variation and the fused lasso replace discrete partition search with convex regularization \citep{rudin1992rof,tibshirani2005fused}; the group fused lasso encourages aligned changes across several sequences \citep{bleakley2011groupfused}. Circular binary segmentation is a standard comparator for array-CGH data \citep{olshen2004cbs}. These methods are important computational and empirical references, but their estimands are penalized or testing-based rather than the posterior distribution over partitions considered here.

\subsection{Irregular designs and design-dependent partition priors}
Generalized product partition models and product partition models with covariates show how covariate information can alter partition probabilities through cohesion functions \citep{park2010gppm,muller2011ppmx}. BayesBreak specializes this idea to ordered observations. It distinguishes a segment-span cohesion from a candidate-boundary hazard because they encode different prior assumptions. Both factors can be included in the ordered-partition probability when they remain local to a segment or transition. Neither is treated as a likelihood weight.

\subsection{Hierarchical and multi-sequence changepoint models}
Hierarchical Bayesian changepoint analysis has a long history \citep{carlin1992hierBayesCP}. BASIC models the tendency of changepoints to co-occur across multiple sequences through a hierarchical empirical-Bayes prior and provides posterior sampling and optimization algorithms \citep{fan2017basic}. Joint random partition models provide a Bayesian nonparametric construction for dependent partitions across several ordered processes \citep{quinlan2024jrpm}. BayesBreak studies a more restrictive setting for exact finite computation: aligned sequences are conditionally independent given a common partition, so their segment marginal likelihoods multiply. Known groups apply this calculation separately by group. The finite-candidate consistency theorem in Section~\ref{sec:extensions} concerns this common-partition model and does not imply consistency for arbitrary hierarchical dependence.

\subsection{Latent groups and common structural changes}
Model-based clustering of time-dependent observations according to common structural changes is an active neighboring literature \citep{corradin2026commonstructural}. The BayesBreak latent-group construction is narrower than a full probability mixture over all partitions. It optimizes a finite objective of the form $\sum_s\log\sum_g\pi_g S_s(\tau_g)$, where $S_s(\tau_g)>0$ is the specified subject--template score. A Jensen minorization yields auxiliary weights and responsibility-weighted max-sum updates. Because the $S_s$ need not be normalized densities over the data space, the algorithm is described as minorization--maximization for a finite latent-group objective, not as maximum likelihood for an identifiable finite mixture \citep{dempster1977em,teicher1963identifiability}.

\subsection{Nonconjugate segment calculations}
When a segment integral is unavailable in closed form, deterministic approximations such as Laplace approximation, the Jaakkola--Jordan variational bound, expectation propagation, P\'olya--Gamma augmentation, or numerical quadrature may be used \citep{jaakkola2000logisticvb,minka2001ep,polson2013pg}. The ordered-partition recursion is unchanged after an approximate segment marginal-likelihood table is supplied, but its output is then exact only for that table. Section~\ref{sec:approx-predict} therefore proves a conditional error-propagation result: a uniform bound on reachable segment log marginal likelihoods yields bounds on partition marginal likelihoods and posterior odds. The theorem does not supply the segmentwise bound for a particular approximation routine.

\subsection{Position of the present work}
The present contribution is a generalized hierarchical Bayesian segmentation model whose observation-family dependence is confined to segment marginal likelihoods and posterior moments. The paper derives the common posterior recursions, the separate MAP recursion, the irregular-design prior factors, exact shared-boundary pooling, the finite latent-group optimization, and the conditional approximation bounds in one notation. Its empirical claims concern the executed simulations and case studies reported here; it does not claim a universal complexity improvement or superiority over every established changepoint method.
